%%%%%%%%%%%%%%%%%%%%%%%%%%%%%%%%%%%%%%%%%%%%%%%%%%%%%%%%%%%%%%%%%%
%% Toto je inkludovaný LaTeXový soubor `ca_conf.tex'.           %%
%% ------------------------------------------------------------ %%
%% Verze 2022-05-23                                             %%
%% Vyvíjí & udržuje <stanislav.hledik@physics.slu.cz>           %%
%%%%%%%%%%%%%%%%%%%%%%%%%%%%%%%%%%%%%%%%%%%%%%%%%%%%%%%%%%%%%%%%%%

\chapter{Konfigurační soubor}\label{config}%%%%%%%%%%%%%%%%%%%%%%%%%%%%%%%%

Balíček \pkg{\ipotname} volitelně používá konfigurační soubor \pgm{\ipotname
  .cfg}. Je-li takto pojmenovaný soubor přítomen v~pracovním adresáři, bude
automaticky nalezen a načten, o~čemž bude uživatel informován hláškou

\medskip

{\small\ttfamily
\noindent
** \ipotname : user config file loaded:\\
(\ipotname .cfg)
\par}

\medskip

\noindent
O~nepřítomnosti konfiguračního souboru bude uživatel informován takto:

\medskip

{\small\ttfamily
\noindent
Class \ipotname{} Warning: `\ipotname .cfg' not found;\\
create one or use preamble to place your stuff.
\par}

\medskip

\noindent
Do \pgm{\ipotname .cfg} si uživatel může uložit vlastní definice a nastavení,
aniž by jimi \uv{plevelil} preambuli hlavního souboru. Prvních 46~řádků
načtených z~aktuálního~\pgm{\ipotname .cfg} pomocí balíčku
\pkg{fancyvrb}~\cite{Vos:LIVE:CTAN:fancyvrb} vypadá takto:

\VerbatimInput[frame=none,firstline=19,lastline=64,fontsize=\footnotesize]%
  {\ipotname .cfg}
\noindent$\langle$\emph{zbývající řádky vynechány}$\rangle$

\medskip

V~první sekci \emph{PDF document information} jsou nastaveny údaje pro
vlastnosti dokumentu zobrazované Adobe Acrobat
Readerem~\cite{Adobe:LIVE::AcrobatReader} i~jinými
prohlížeči~\designation{pdf}, např.~M$\mu$PDF~\cite{Art:LIVE::MuPDF}. Další
sekce \emph{Page positioning} je připravena pro jemné doladění umístění
sazebního obrazce popsané v~Dodatku~\ref{mirror}. Sekce \emph{TOC and
  sectioning} je připravena na změnu hloubky Obsahu a číslování sekcí. Sekce
\emph{Repeatedly used URLs} a \emph{Base URLs for arXiv and doi} definuje
opakovaně používaná URL a bázová URL pro \designation{arXiv} a
\designation{doi}. Následuje sekce \emph{General text macros and
  abbreviations} s~obecně použitelnými makry.

Soubor \pgm{\ipotname .cfg} dodávaný v~balíčku \pkg{\ipotname} je mnohem
rozsáhlejší, další materiál lze nalézt přímo v~něm. Předpokládá se, že
uživatel si jej přizpůsobí pro své vlastní potřeby. Zejména je záhodno změnit
údaje v~sekci \emph{PDF document information}. Tyto údaje se uvádějí zvlášť
(místo aby byly vzaty z~údajů pro titulní stránky) z~toho důvodu, že vyžadují absenci diakritiky,
která je ve vlastnostech dokumentu chybně zobrazována. Zakomentujete-li
některý z~příkazů \pgm{\textbackslash pdfinfo*}, zůstane odpovídající políčko
ve vlastnostech dokumentu prázdné.



\endinput

%%
%% End of file `ca_conf.tex'.
